\fsection{Ej 4 pág. 150}
\fsubsection
[\textcolor{waveBlue2}{Investiga}]
{\textcolor{waveBlue2}{\boldit{Investiga}}}
{
	Entra en el portal de la \textsc{Asociación para la Sostenibilidad y el Progreso de las
		Sociedades} (\textless\,\url{https://sostenibilidadyprogreso.org}\,\textgreater) y busca
	información sobre su función. Elabora un pequeño dosier sobre sus objetivos y para qué
	fue creada la asociación.
}

\begin{solution}

	\textbf{DOSIER: Asociación para la Sostenibilidad y el Progreso de las Sociedades (ASYPS)}

	\medskip

	\textbf{¿Qué es ASYPS?}

	ASYPS es una organización de la sociedad civil creada para actuar como instrumento de análisis,
	investigación y foro de debate en materia de sostenibilidad y desarrollo sostenible. Funciona como
	plataforma de comunicación y centro de observación permanente sobre los procesos de sostenibilidad
	y progreso de las sociedades, desde una visión global e integrada de las interacciones ambientales,
	económicas, sociales, culturales e institucionales.

	\medskip

	\textbf{¿Por qué fue creada?}

	Nació ante la necesidad de contar con capacidades independientes y competentes desde la sociedad
	civil organizada para hacer frente al Cambio Global y a la crisis sistémica actual. Su fundación
	surgió siguiendo la estela del Observatorio de la Sostenibilidad en España (OSE), con el propósito
	de dar continuidad a sus planteamientos sobre sostenibilidad integral.

	\medskip

	\textbf{Objetivos principales:}

	\begin{itemize}
		\item Facilitar información relevante, rigurosa e independiente sobre sostenibilidad y progreso
		      de las sociedades.
		\item Fomentar el debate social para mejorar la toma de decisiones de administraciones, agentes
		      económicos y sociales, y ciudadanía en general.
		\item Impulsar respuestas y soluciones ante el Cambio Global y favorecer la transición
		      socioecológica hacia nuevos paradigmas de progreso sostenible.
		\item Desarrollar proyectos sociales y empresariales que sirvan para transformar los modos de
		      producción, consumo, educación y comportamiento con criterios solidarios, responsables y
		      sostenibles.
		\item Ofrecer espacios para el debate creativo y el impulso de proyectos innovadores que
		      promuevan la sostenibilidad y el progreso social.
	\end{itemize}

	\medskip

	\textbf{Actividades destacadas:}

	ASYPS organiza foros, ciclos de conferencias (entre ellos, uno anual sobre Economía Circular),
	jornadas temáticas y publica el \textit{Boletín Observatorio ASYPS}, un boletín digital de
	difusión de temas de actualidad en sostenibilidad, así como tribunas de opinión y documentos
	de trabajo.

	\medskip

	\textbf{Principios de actuación:}

	Sostenibilidad, ética ecológica, diversidad cultural, cohesión social, equidad y solidaridad
	intra e intergeneracional.

\end{solution}
